%% @q {"answer": "B", "figures": [], "hash": "4c85d49244168613", "key": "demo/示例卷#1", "meta": {"book": "demo", "source_label": "示例卷", "source_no": 1}, "options": [{"label": "A", "text": "$\\{1\\}$"}, {"label": "B", "text": "$\\{2\\}$"}, {"label": "C", "text": "$\\{3\\}$"}, {"label": "D", "text": "$\\{1,2,3\\}$"}], "points": [], "solution": "两个集合的公共元素只有 $2$，故 $A\\cap B=\\{2\\}$。故选 B。", "stem": "已知集合 $A=\\{1,2\\}$，$B=\\{2,3\\}$，则 $A\\cap B=$\\paren[B]", "type": "single_choice"}
\begin{question}
已知集合 $A=\{1,2\}$，$B=\{2,3\}$，则 $A\cap B=$\paren[B]
\begin{choices}
  \item $\{1\}$
  \item $\{2\}$
  \item $\{3\}$
  \item $\{1,2,3\}$
\end{choices}
\begin{solution}
两个集合的公共元素只有 $2$，故 $A\cap B=\{2\}$。故选 B。
\end{solution}
\end{question}


%% @q {"answer": "ABC", "figures": [], "hash": "7c34cb334fdfec24", "key": "demo/示例卷#9", "meta": {"book": "demo", "source_label": "示例卷", "source_no": 9}, "options": [{"label": "A", "text": "$ab\\leqslant 1$"}, {"label": "B", "text": "$a^{2}+b^{2}\\geqslant 2$"}, {"label": "C", "text": "$\\dfrac{1}{a}+\\dfrac{1}{b}\\geqslant 2$"}, {"label": "D", "text": "$a^{3}+b^{3}\\leqslant 2$"}], "points": [], "solution": "由 $a+b=2$ 与均值不等式：$ab\\leqslant\\left(\\dfrac{a+b}{2}\\right)^{2}=1$，A 正确；$a^{2}+b^{2}\\geqslant\\dfrac{(a+b)^{2}}{2}=2$，B 正确；$\\dfrac{1}{a}+\\dfrac{1}{b}=\\dfrac{a+b}{ab}=\\dfrac{2}{ab}\\geqslant 2$，C 正确；取 $a=2$、$b\\to 0^{+}$ 时 $a^{3}+b^{3}\\to 8>2$，D 错误。故选 ABC。", "stem": "已知 $a>0$，$b>0$，且 $a+b=2$，则下列结论正确的是\\paren[ABC]", "type": "multi_choice"}
\begin{question}
已知 $a>0$，$b>0$，且 $a+b=2$，则下列结论正确的是\paren[ABC]
\begin{choices}
  \item $ab\leqslant 1$
  \item $a^{2}+b^{2}\geqslant 2$
  \item $\dfrac{1}{a}+\dfrac{1}{b}\geqslant 2$
  \item $a^{3}+b^{3}\leqslant 2$
\end{choices}
\begin{solution}
由 $a+b=2$ 与均值不等式：$ab\leqslant\left(\dfrac{a+b}{2}\right)^{2}=1$，A 正确；$a^{2}+b^{2}\geqslant\dfrac{(a+b)^{2}}{2}=2$，B 正确；$\dfrac{1}{a}+\dfrac{1}{b}=\dfrac{a+b}{ab}=\dfrac{2}{ab}\geqslant 2$，C 正确；取 $a=2$、$b\to 0^{+}$ 时 $a^{3}+b^{3}\to 8>2$，D 错误。故选 ABC。
\end{solution}
\end{question}


%% @q {"answer": "-1", "figures": [], "hash": "29ccdcf0925da8a6", "key": "demo/示例卷#12", "meta": {"book": "demo", "source_label": "示例卷", "source_no": 12}, "options": [], "points": [], "solution": "$f(x)=(x-1)^{2}-1$，其图象开口向上、对称轴为 $x=1$，而 $1\\in[0,3]$，故最小值为 $f(1)=-1$。", "stem": "已知函数 $f(x)=x^{2}-2x$，则 $f(x)$ 在 $[0,3]$ 上的最小值为\\fillin[-1]", "type": "fill_in_blank"}
\begin{question}
已知函数 $f(x)=x^{2}-2x$，则 $f(x)$ 在 $[0,3]$ 上的最小值为\fillin[-1]
\begin{solution}
$f(x)=(x-1)^{2}-1$，其图象开口向上、对称轴为 $x=1$，而 $1\in[0,3]$，故最小值为 $f(1)=-1$。
\end{solution}
\end{question}


%% @q {"answer": "见解析", "figures": [], "hash": "38122914acf41c36", "key": "demo/示例卷#15", "meta": {"book": "demo", "figures": ["7e044a2c74964f15.png"], "source_label": "示例卷", "source_no": 15}, "options": [], "points": [], "solution": "(1) 令 $y=0$，得 $x^{2}-2x-3=0$，解得 $x=-1$ 或 $x=3$，故 $A(-1,0)$、$B(3,0)$。\n\n(2) 令 $x=0$，得 $y=-3$，故 $C(0,-3)$。于是\n\\[\nS_{\\triangle ABC}=\\dfrac{1}{2}|AB|\\cdot|OC|=\\dfrac{1}{2}\\times 4\\times 3=6 .\n\\]", "stem": "如图，抛物线 $y=x^{2}-2x-3$ 与 $x$ 轴交于 $A$、$B$ 两点（$A$ 在 $B$ 的左侧），与 $y$ 轴交于点 $C$。\n\n(1) 求 $A$、$B$ 两点的坐标；\n\n(2) 求 $\\triangle ABC$ 的面积。\n\n\\includegraphics[width=0.5\\linewidth]{7e044a2c74964f15.png}", "type": "detailed_answer"}
\begin{problem}
如图，抛物线 $y=x^{2}-2x-3$ 与 $x$ 轴交于 $A$、$B$ 两点（$A$ 在 $B$ 的左侧），与 $y$ 轴交于点 $C$。

(1) 求 $A$、$B$ 两点的坐标；

(2) 求 $\triangle ABC$ 的面积。

\includegraphics[width=0.5\linewidth]{7e044a2c74964f15.png}
\begin{solution}
\textbf{答案：}见解析
(1) 令 $y=0$，得 $x^{2}-2x-3=0$，解得 $x=-1$ 或 $x=3$，故 $A(-1,0)$、$B(3,0)$。

(2) 令 $x=0$，得 $y=-3$，故 $C(0,-3)$。于是
\[
S_{\triangle ABC}=\dfrac{1}{2}|AB|\cdot|OC|=\dfrac{1}{2}\times 4\times 3=6 .
\]
\end{solution}
\end{problem}

